\subsection{Exact five-primary support matching in the original recursion}

\begin{result}{E.211}{Proved}
{The first finite system is the only obstruction to five-primary class-two support matching in this original subfamily}
Fix legal first primes $q_0,r_0$ in the original factorial recursion
E.102, with the minimal five-adic shells of E.199. The following
conditions are equivalent.
\begin{enumerate}
\item The affine system E.202.1 is consistent.
\item There is a complete continuation of the unchanged original
recursion, retaining the minimal shells at every stage, whose unique
global integral pro-five class-two section class agrees with its
ordinary local point section at every support prime.
\end{enumerate}
The assertion concerns this one coefficient prime. In particular it
does not supply a legal first prime pair satisfying the first
condition, a match at the coefficient place five, or matching at
any higher nilpotent class.

More precisely, after an actual later pair $q_j,r_j$ has been chosen,
where $j\ge1$, put
\[
 m=n_j=(j+4)!,\qquad e=v_5(n_{j+1}),\qquad R=\mathbb Z/5^e.
\]
Choose these primes with the additional unit conditions of E.205:
$\alpha=a(F_{q_j})\in\mathbb Z_5^\times$ and
$\beta=b(F_{r_j})\in\mathbb Z_5^\times$.
Let $z_q,z_r$ be the residue-one basis of
$H^1(G_{\mathbb Q,\{5,q_j,r_j,\infty\}},R(2))$.
If all previously corrected old roots are thereafter chosen to be
principal units, the exact changes of the two current mismatch
values, as the next complementary roots vary, are
\[
 \begin{pmatrix}\Delta D_{A,j}\\\Delta D_{B,j}\end{pmatrix}
 =\begin{pmatrix}
 m z_q(F_{q_j})-q_j\alpha&-m z_r(F_{q_j})\\
 m z_q(F_{r_j})&\beta-m z_r(F_{r_j})
 \end{pmatrix}
 \begin{pmatrix}\Delta u_j\\\Delta v_j\end{pmatrix}
 \quad\text{in }R^2.
 \tag{E.211.1}
\]
Both primitive increments range independently over all of $R$.
The matrix is invertible, since its reduction modulo five is
$\operatorname{diag}(-\alpha,\beta)$.
Thus an exact prescribed current pair of values in $R^2$ can be
obtained without changing any earlier support match. In particular
both current values can be made zero.
\end{result}

\begin{proof}
\emph{The local target and the finite residue basis.}
Write $q=q_j,r=r_j$. The minimal shells give
\[
 v_5(q-1)=v_5(r-1)=v_5(q^2-1)=v_5(r^2-1)=e.
 \tag{E.211.2}
\]
The exact ordinary local class-two comparison of B.24 says that
the $A$-support point is matched if and only if
\[
 D_{A,j}=h(F_q)-q\alpha u_j\in5^e\mathbb Z_5,
 \qquad b(F_q)=(q-1)u_j.
 \tag{E.211.3}
\]
At the $B$-support point its condition is
\[
 h(F_r)+\beta v_j-\eta(F_r)\in5^e\mathbb Z_5,
 \qquad a(F_r)=(r-1)v_j.
\]
Here $\eta(F_r)=(r^2-1)/24$ belongs to $5^e\mathbb Z_5$.
Consequently this is exactly
$D_{B,j}=h(F_r)+\beta v_j=0$ in $R$, with the same
normalization used in E.199 and E.205. These are actual local
point matches, not merely the weaker class-three residue tests.
The local weight-one invariants vanish integrally, so their
meaning is unaffected by unrecorded lower conjugations.

For $S_5=\{5,\infty\}$ the fixed-source calculations in E.202
(or E.209) give
\[
 H^1(G_{\mathbb Q,S_5},R(2))=
 H^2(G_{\mathbb Q,S_5},R(2))=0.
 \tag{E.211.4}
\]
For clarity the first equality follows from
$H^1(G_{\mathbb Q,S_5},\mathbb Z_5(2))=0$ and the vanishing
of integral $H^2$, by its multiplication-by-$5^e$ exact sequence.
The second is the literal finite assertion of E.209, also
obtainable by the coefficient calculation in E.202.
Finite purity and E.211.2 therefore identify
\[
 H^1(G_{\mathbb Q,S_5\cup\{q,r\}},R(2))
 \xrightarrow{\sim}
 H^0(\mathbb F_q,R(1))\oplus H^0(\mathbb F_r,R(1))=R^2.
 \tag{E.211.5}
\]
This defines the two basis classes in the statement.
Every central coboundary has Frobenius value
$(p^2-1)c=0$ in $R$ at $p=q,r$, so the displayed evaluations
are well-defined on cohomology classes in the fixed frames.

\emph{What is fixed before the next primes are chosen.}
Take an actual prefix through $q_j,r_j$. All older corrected
support places are required to use roots congruent to one modulo
that place at every later stage. This is a permitted choice in
E.102: initially the complementary derivative is a unit, and when
the old place eventually divides the factorial, the original
stronger error bound gives the required principal roots by Newton
iteration. These requirements retain the original current-prime
congruence and every radical-splitting condition.

At the next step, whose exponent is $n_{j+1}$, choose any one
legal pair of old roots at each of $q,r$. At $q$ allow its
complementary $Y$-root to be multiplied by an arbitrary element
of $\mu_{5^e}(\mathbb Q_q)$; at $r$ allow its complementary
$X$-root to be multiplied independently by an element of
$\mu_{5^e}(\mathbb Q_r)$. These roots of unity exist by
E.211.2 and preserve the exact old equations. The original
choice gave $q_j,r_j>n_{j+2}$, so neither place divides the
next factorial $n_{j+2}$. No additional congruence requiring
these particular roots to be one modulo that place is violated.
After this next step, require principal roots at $q,r$ too.

Every such pair of root choices has a complete legal ordinary
continuation, by E.102. The extra unit Frobenius conditions
for future new pairs may be imposed simultaneously, by the
exact field-intersection argument of E.205. That argument keeps
the original full CRT field and radical splitting: the
$A_0^{1/5}$ line has private ramification at three, whereas
the complementary radical's degree-five line is unramified
there; the corresponding $B_0^{1/5}$ argument uses two.
Thus these additional conditions do not restrict the two
current root-of-unity choices.

We can consequently use any one of these complete continuations
as a provisional reference for the present finite calculation.
Nothing asserts that its later matches have already been
corrected. The following comparison shows that the result of
the current finite calculation is independent of its arbitrary
later primes, within the stated principal-root requirements.

\emph{Literal cochains and every inertia difference.}
Choose the global integral Kummer cocycles as the uniformly
convergent sums of the fixed initial terms and the individual
prime Kummer cocycles multiplied by their factorial exponents.
For any two completions just described, the abelian cocycles
are literally the same modulo $5^e$: the actual prefix through
$q_j,r_j$ is fixed and every later exponent is divisible by
$n_{j+1}$. Thus the difference $h'-h$ of their integral
class-two cochains reduces to an $R(2)$-valued one-cocycle.
The integral class-two class at five is unique, since
$H^1(G_{\mathbb Q},\mathbb Z_5(2))=0$. Its central coboundary
ambiguity has zero value at the two current Frobenii modulo
$5^e$ by E.211.2.

We verify that this difference is unramified outside
$\{5,q,r,\infty\}$. At a previously corrected support prime
$p=q_i$ or $r_i$, $i<j$, every ratio of subsequent prime
factors is a principal unit at $p$. Its local Kummer cocycle
has an integral primitive. The primitive change is a convergent
sum of these integral values multiplied by exponents divisible
by $n_{j+1}$, hence belongs to $5^e\mathbb Z_5$.
The exact class-two tame relation
\[
             h(T_{q_i})=n_i u_i,
             \qquad h(T_{r_i})=-n_i v_i
 \tag{E.211.6}
\]
then makes the central inertia difference zero modulo $5^e$.
At the fixed primes two and three the same conclusion follows
either directly from their integral unit primitives or from
the tame relation with $p-1$ a five-adic unit and identical
abelian cocycles modulo $5^e$.

Every later support prime has exponent divisible by
$n_{j+1}$; its integral divided complementary primitive gives
$h(T)=0$ modulo $5^e$ individually on each branch. A prime
which has zero valuation in a branch has $h(T)=0$ integrally,
by the torsion-free class-two tame relation. These statements
also cover a prime appearing in only one branch or at different
later indices in two branches. Wild and prime-to-five tame
inertia vanish for the actual pro-five cocycles away from five.
These are all places outside the two current primes and five.
Therefore $h'-h$ descends to the source of E.211.5.

At the two current places its residue vector is exactly
\[
                  (m\,\Delta u_j,-m\,\Delta v_j).
 \tag{E.211.7}
\]
Equation E.211.5 gives
\[
 [h'-h]=m\,\Delta u_j z_q-m\,\Delta v_j z_r
       \quad\text{in }H^1(G_{\mathbb Q,S_5\cup\{q,r\}},R(2)).
 \tag{E.211.8}
\]
In particular if the two primitive increments vanish in $R$,
the Frobenius changes vanish there too. This is the promised
independence from the unrestricted later prime choices.

\emph{The exact finite affine solve.}
The Kummer cocycle of a chosen primitive $5^e$th root of
unity at $q$ has divided Frobenius primitive $1/5^e$, up to
an integral primitive. Multiplying its exponent by
$n_{j+1}$ changes $u_j$ by
\[
                 (n_{j+1}/5^e)i\pmod{5^e},\qquad i\in R.
 \tag{E.211.9}
\]
The coefficient is a five-adic unit. The finite CRT precision
retains the selected root of unity in the residue field;
the remaining factor is a principal unit and contributes
$n_{j+1}$ times an integral primitive, which is zero in $R$.
The identical calculation at $r$ changes $v_j$ through all
of $R$, independently. All later principal-root factors
have the same zero contribution in $R$.

The abelian Frobenius values $\alpha,\beta$ are common
modulo $5^e$, since the abelian cochains agree literally.
Taking Frobenius values in E.211.8 and substituting in
E.211.3 gives exactly E.211.1. Since $j\ge1$, $5\mid m$;
since $q\equiv1\pmod5$ and $\alpha,\beta$ are units, the
displayed matrix has unit determinant in $R$. Every desired
current mismatch vector therefore has one and only one
primitive-increment parameter relative to the chosen reference.
Choose the parameter which makes both values zero.

This operation retains every earlier match. For $i<j$,
\[
                   e_i=v_5(n_{i+1})\le v_5(n_j).
 \tag{E.211.10}
\]
At that earlier precision the changes at the two current
inertias in E.211.7 vanish, while all other inertia changes
were just shown to vanish. The difference is unramified
outside five, so E.211.4 makes it a central coboundary at
that precision. At the earlier Frobenius this coboundary
is zero, because $p^2-1$ is divisible by $5^{e_i}$.
Its divided complementary primitive is unchanged modulo
$5^{e_i}$ by the principal-root requirement. Hence the
earlier mismatch remains zero. The same argument shows
that subsequent corrections retain the newly obtained
match as well.

\emph{The infinite original continuation and its exact scope.}
Suppose E.202.1 is consistent. Choose its successful stage-one
root parameter, and choose $q_1,r_1$ with the unit conditions
of E.205. The first shell has $e_0=1$, so its two finite
conditions already give the full local class-two matches
at $q_0,r_0$. Keep principal roots at those places thereafter.
Apply the just-proved finite solve successively for $j=1,2,\ldots$,
choosing each next prime pair with the same unit conditions.
Each step is a permitted original CRT/radical step, and
E.211.10 preserves all earlier matches. The original proof
E.102 gives the resulting ordinary infinite-support completed
pair and its exact ordinary local values.

At two and three the local integral group
$H^1(G_{\mathbb Q_p},\mathbb Z_5(2))$
is $\mathbb Z_5/(p^2-1)=0$, so the global class-two section
also agrees with the ordinary point class there. These
and the constructed pairs are the entire support.
Conversely, any original continuation with those class-two
matches passes the first two finite tests; the all-continuation
exhaustion in E.202 then forces consistency of E.202.1.
This proves the equivalence.

Boundary matching also implies the strict class-three
convergence condition by B.24, and E.209 removes its
fixed-source obstruction. No claim about choosing that
class-three lift to match every support point follows:
its separate matching problem is still B.26, with the
parameters of E.210. This result strengthens E.205 at the
class-two boundary gate. E.213 supplies a first pair satisfying
its finite condition; the subsequent tower conditions remain
separate.
\end{proof}
